			\TemplateSection[K0E510E020E990E2C0702CCF4C4CF8325010112010100]{Math}
				\TemplateSubsection[KCF4214C7BB0DCCD931CB773AD139310727CBDC2BCFF837C90B49CCF4C4072A0101010100]{行列式求值(任意模数)}
					\TemplateCode{cpp}{src/Math/行列式求值(任意模数）.cpp}
				\TemplateSubsection[KC43B19CD2F79CF165ED0560D0101010100]{高斯消元}
					\TemplateCode{cpp}{src/Math/高斯消元.cpp}
				\TemplateSubsection[K0E480E7C0E700E2507020E480E7C0E700E2507020E7C07270D1A072A0702C1B06D072F07020E090E020E8A0E8A0E210E990E99010112020202021202020202120202020202020212010100]{Long Long $O(1)$ 乘, Barrett}
					\TemplateCode{cpp}{src/Miscellany/LLFPM.cpp}
				\TemplateSubsection[K0E210EA60E250E0A0E1A072F0702C99834D0560D0101010100]{exgcd, 逆元}
					\TemplateNote{src/Math/exgcd.tex}
					\TemplateCode{cpp}{src/Math/exgcd.cpp}
					递推逆元: $\operatorname{inv}(i) \equiv (P - P / i) \cdot \operatorname{inv}(P \bmod i)$
				\TemplateSubsection[KCE1A13D04143C3B610C1B0AC0101010100]{同余方程}
					$ax \equiv b \bmod c$
					\TemplateCode{cpp}{src/Math/同余方程.cpp}
				\TemplateSubsection[K0E0A0E8A0E990702D1490DC4C925CCCB40D04143C31F28C77B520101121212010100]{CRT 中国剩余定理}
					\TemplateCode{cpp}{src/Math/CRT_lbn.cpp}
				\TemplateSubsection[K0E510E320E480E480E210E8A07020E8A0E020E090E320E70072F07020E7E0E7C0E480E480E020E8A0E1A07020E8A0E2C0E7C01011202020202020212020202020202120202020202020212010100]{Miller Rabin, Pollard Rho}
					\TemplateCode{cpp}{src/Math/Miller Rabin And Pollard Rho.cpp}
				\TemplateSubsection[KCF062BCF481CC5AFE20101010100]{线性基}
					\TemplateCode{cpp}{src/Math/线性基.cpp}
				\TemplateSubsection[KC7120DD0F119C8080DC6740DCD2F790101010100]{扩展卢卡斯}
					\TemplateCode{cpp}{src/Math/扩展卢卡斯.cpp}
				\TemplateSubsection[KC78F10CF6BD3C71513C44731C73D13CBE80DC15376D139310101010100]{连续拉格朗日插值}
					\TemplateCode{cpp}{src/Math/连续拉格朗日插值.cpp}
				\TemplateSubsection[KC5FC0DC1B06DCB8913C90B490101010100]{阶乘取模}
					\TemplateCode{cpp}{src/Math/Factorial Mod.cpp}
				\TemplateSubsection[KD1416DCCF4C4C44E0DCCF4C4C5BE16CCF4C40101010100]{质数个数计数}
					时间复杂度$O(n^{\frac{2}{3}})$, 可以处理大概 $10^{13}$ 的范围。
					\TemplateCode{cpp}{src/Math/质数个数.cpp}
				\TemplateSubsection[KC7601FCA2619C5BB0DC77B1FC2CB2E0702D13916CF062BCEF713C44731C2F419CE1E10C5BE160101010100]{类欧几里得 直线下格点统计}
					\TemplateCode{cpp}{src/Math/直线下格点统计.cpp}
				\TemplateSubsection[KCE8D0DC97D0DCA2619C5BB0DC77B1FC2CB640101010100]{万能欧几里德}
					\TemplateCode{cpp}{src/zjj/euclid.cpp}
				\TemplateSubsection[KCAD113C3B610CCCB40D041430101010100]{平方剩余}
					\TemplateCode{cpp}{src/Math/平方剩余.cpp}
				% \subsection{BSGS 离散对数}
				% \inputminted{cpp}{src/Math/BSGS.cpp}
				\TemplateSubsection[KCF062BCF481CCE1A13D04143C1080DC2DA1CCCD9310101010100]{线性同余不等式}
					\TemplateCode{cpp}{src/Math/线性同余不等式.cpp}
				\TemplateSubsection[K0E510E320E7007020D1C0D220702CC790D0101010100]{min\_25 筛}
					\TemplateNote{src/Math/min25_notes.tex}
					\TemplateCode{cpp}{src/Math/min25_seive.cpp}
				% \inputminted{cpp}{src/Math/TEES.cpp}
				\TemplateSubsection[KD0565BC454130101010100]{原根}
				% \inputminted{cpp}{src/Math/原根.cpp
					\TemplateNote{src/tbr/阶和原根.tex}
				\TemplateSubsection[KC05E10D0861CCF062BC64E16C5AFCA0101010100]{半在线卷积}
					\TemplateCode{cpp}{src/Math/半在线卷积.cpp}
				\TemplateSubsection[KD0861CCF062BC64E16C5AFCA0101010100]{在线卷积}
					\TemplateCode{cpp}{src/Math/在线卷积.cpp}
				\TemplateSubsection[KC3680DCF1319CCD9310101010100]{多项式}
					\TemplateCode{cpp}{src/Math/poly.cpp}
				\TemplateSubsection[K0E230E230E990101121212010100]{FFT}
					\TemplateCode{cpp}{src/Math/FFT.cpp}
				\TemplateSubsection[K0E700E990E990727C29E0DC90B49CCF4C4072A0101121212010100]{NTT(大模数)}
					\TemplateCode{cpp}{src/Math/NTT_big.cpp}
				\TemplateSubsection[K0E510E990E990702CBDC2BCFF837C90B49CCF4C4C64E16C5AFCA0101121212010100]{MTT 任意模数卷积}
					\TemplateCode{cpp}{src/Math/MTT.cpp}
				\TemplateSubsection[KC3680DCF1319CCD931D07416CD64250101010100]{多项式运算}
					\TemplateSubsubsection{多项式求逆\ 开根\ ln exp}
						\TemplateCode{cpp}{src/Math/多项式运算.cpp}
					\TemplateSubsubsection{多项式除法\ 取模}
						需要抄求逆。
						\TemplateCode{cpp}{src/Math/多项式取模.cpp}
					\TemplateSubsubsection{多点求值}
						需要抄取模。
						\TemplateCode{cpp}{src/Math/多点求值.cpp}
					\TemplateSubsubsection{插值}
						\TemplateNote{src/Math/插值.tex}
					% \subsubsection*{快速插值}
					% 	\input{src/Math/快速插值.tex}
				\TemplateSubsection[KCF062BCF481CC2EA76CE490D0101010100]{线性递推}
					\TemplateSubsubsectionStar{$O(k^2 \log n)$}
						\TemplateCode{cpp}{src/Math/线性递推.cpp}
					\TemplateSubsubsectionStar{$O(k \log k \log n)$ - Bostan-Mori}
						\TemplateCode{cpp}{src/Math/快速线性递推-bostan-mori.cpp}
					\TemplateSubsubsectionStar{$O(k \log k \log n)$ - 多项式取模}
						需要抄前面的多项式取模。预处理 $O(k \log k \log n)$，固定 $n$ 和系数只改变初始值的话，询问一次 $O(k)$。注意只询问一次的话不如 Bostan-Mori 快。
						\TemplateCode{cpp}{src/Math/快速线性递推-多项式取模.cpp}
				\TemplateSubsection[K0E090E210E8A0E480E210E360E020E510E7E0E510E020E910E910E210EA70702D1E62BCF1E0DC3680DCF1319CCD9310101120202020202020202120101FFF6821200]{Berlekamp-Massey 最小多项式}
					\TemplateNote{src/Math/Berlekamp-Massey.tex}
				\TemplateSubsection[K0E230EA40E990101121212010100]{FWT}
					\TemplateCode{cpp}{src/Math/FWT.cpp}
					\TemplateCode{cpp}{src/Math/myFWT.cpp}
				% \subsection{FFT NTT FWT 多项式操作}
				% 	\inputminted{cpp}{src/Math/FFT NTT FWT.cpp}
				\TemplateSubsection[K0E360702C61516D1416707020E230EA40E9901011202020202121212010100]{K 进制 FWT}
					\TemplateCode{cpp}{src/tbr/fwt.cpp}
				% \input{src/Math/多项式牛顿法.tex}
				% \inputminted{cpp}{src/Math/Polynomial.cpp}
					% \input{src/Math/单位根反演.tex}
				%\input{src/Math/多点求值与快速插值.tex}
				%\subsection{多项式插值}
				%\input{cpp}{src/Math/插值.tex}
				\TemplateSubsection[K0E910E320E510E7E0E480E210EA60702C2AE3AC22F19CF4216010112010100]{Simplex 单纯形}
					\TemplateCode{cpp}{src/Math/Simplex.cpp}
				\TemplateSubsection[KC43B19CD2F79CF165ED0560DD1E62BCF1E0DC3B346CCF4C4C6071F0101010100]{高斯消元最小范数解}
					\TemplateCode{cpp}{src/Geometry/minnorm_gauss.cpp}
				%\newpage
				\TemplateSubsection[K0E7E0E210E480E480702C3B610C1B0AC010112010100]{Pell 方程}
					\TemplateCode{cpp}{src/Math/Pell方程.cpp}
				\TemplateSubsection[KC6071FCFE60DD0560DCC3E0DC24516C3B610C1B0AC0101010100]{解一元三次方程}
					\TemplateCode{cpp}{src/Math/解一元三次方程.cpp}
				\TemplateSubsection[KD1B510CCD9AFD0101207020E910E320E510E7E0E910E7C0E7001010202020212010100]{自适应 Simpson}
					\TemplateCode{cpp}{src/Math/Simpson.cpp}
				%\subsection{Schreier-Sims 求排列群的阶}
				%\inputminted{cpp}{src/Math/Schreier-Sims.cpp}
			\newpage
